%% notes-tex/031-unified-representation/sections/2b-what-a-record-is.tex

\section{What a record is and what is done to it\secstatus{tent}}
\label{sec:record}

Before any argument about pooling, this section fixes what the data actually looks like and
what the model is asked to do with it, because the rest of the document is unreadable without
both. Table~\ref{tab:vocab} defines every term that appears in a later figure label or table
header.

%% SOURCE: VOCABULARY in experiments/031-env-chemgen-inhibitor-tolerance/scripts/notes_tex_representation_tables.py
%% (curated definitions; the measured quantities they name are produced by the scripts cited
%% at each table and figure).
\begin{table}[htbp]
\centering
\footnotesize
\caption{Controlled vocabulary, alphabetical. Each term is used in a figure label or a table
header later in the document.}
\label{tab:vocab}
%% GENERATED by experiments/031-env-chemgen-inhibitor-tolerance/scripts/notes_tex_representation_tables.py
%% SOURCE: experiments/031-env-chemgen-inhibitor-tolerance/results/VOCABULARY in notes_tex_representation_tables.py (curated definitions; the measured quantities they name are produced by the scripts cited elsewhere)
%% Do not edit by hand.
\begin{tabular}{p{3.0cm}p{13.4cm}}
\toprule
\textbf{term} & \textbf{meaning} \\
\midrule
calculated logP & The octanol-water partition coefficient estimated from structure by the Crippen method in RDKit. It measures lipophilicity: a negative value is water-preferring, a positive value is fat-preferring. Nothing is measured here; it is computed. \\
ceiling & The highest correlation any model could reach against the noise-free response, given how noisy the measurement is. It is the square root of the reliability. A score of 0.31 against a ceiling of 0.84 means the model captured 37 percent of what the data allows, not 31 percent of the truth. \\
cell & One (gene, compound) pair. Used when counting what two datasets both measured, so that a comparison is between the same gene under the same compound. \\
compound cold-start & An evaluation that holds out every record of one compound at a time, so the model must say something about a molecule it never saw dosed. This is the split that matters for predicting a hydrolysate inhibitor. \\
HIP, HOP & The two assays of a pooled chemogenomic screen. HIP (haploinsufficiency profiling) uses heterozygous diploids, one wild-type copy left. HOP (homozygous profiling) uses homozygous diploids, both copies gone. \\
HOM, HET & The two arms of Hillenmeyer 2008, its homozygous and heterozygous collections, released as separate files with different readouts and handled here as separate datasets. The homozygous arm is the one dropped. \\
MADL score & Median absolute deviation logarithmic score, the per-experiment readout Hoepfner reports: the strain's log abundance ratio minus the median over all strains in the sample, divided by the median absolute deviation of that same set. \\
orientation factor & Plus or minus one per dataset, multiplied into the response so that afterwards a more negative value means a sicker strain everywhere. It is a sign flip on a continuous number, not a threshold and not a class label. \\
reliability & The share of a measurement's variance that is real rather than noise. Estimated either from a served standard error or from the agreement between two independent measurements of the same condition. \\
sensitized host & A strain whose drug-efflux regulators are deleted so that compounds accumulate instead of being pumped out. Vanacloig deletes PDR1, PDR3 and SNQ2, so every one of its genotypes carries those three deletions plus the gene being queried, which is why its records perturb four genes rather than one. \\
\bottomrule
\end{tabular}

\end{table}

\notesec{The task is regression on a continuous number}
It is worth stating plainly, because the sign conventions in the next section invite the
opposite reading. The model predicts a real-valued growth response for a strain under a
compound. Nothing is thresholded into sick and healthy, no class label is formed, and no
cutoff appears anywhere in the pipeline. The reason the document talks about which SIGN means
a sick strain is that the five sources disagree about that sign, and a pooled target has to
resolve the disagreement before it can be regressed on. Resolving it costs nothing: the fix
is a multiplication by plus or minus one, which is a bijection on the real line and discards
no information. Were the target binarized, the compound-level correlations that every score
in this document reports would not be the right way to measure it.

\notesec{One real record from each dataset, carried through}
Table~\ref{tab:worked} takes an actual served record from each of the five and prints it at
every step. Reading a row left to right gives what the strain is, what it was dosed with,
what the source reported, and what the pooled model would be asked to predict.

%% SOURCE: experiments/031-env-chemgen-inhibitor-tolerance/scripts/worked_example_records.py ->
%% results/worked_example_records.csv, rendered by notes_tex_representation_tables.py.
\begin{table}[htbp]
\centering
\footnotesize
\caption{One real served record from each dataset, carried through the two target
transformations. The strain column gives the queried gene with its common name, and for
Vanacloig also the three sensitized-host deletions every one of its genotypes carries.
$o_k$ is the orientation factor and $\tilde y$ the standardized target. The records were
chosen as the strongest response of a gene measured in all five, so the arithmetic is visible
rather than sitting at zero.}
\label{tab:worked}
%% GENERATED by experiments/031-env-chemgen-inhibitor-tolerance/scripts/notes_tex_representation_tables.py
%% SOURCE: experiments/031-env-chemgen-inhibitor-tolerance/results/worked_example_records.csv from worked_example_records.py
%% Do not edit by hand.
\begin{tabular}{p{2.6cm}p{3.0cm}llp{2.3cm}rrr}
\toprule
\textbf{dataset} & \textbf{strain} & \textbf{ploidy} & \textbf{compound} & \textbf{reports} & \textbf{raw} & \textbf{$o_k$} & \textbf{$\tilde y$} \\
\midrule
Vanacloig 2022 & YBR058C (UBP14) + host & hap & crystal violet & log2 ratio & -3.40 & +1 & -4.80 \\
Hillenmeyer HET & YBR058C (YBR058C) & dip & nocodazole & log2 ratio & +1.92 & -1 & -4.30 \\
Hoepfner 2014 & YBR058C (YBR058C) & dip & Epothilon B Deriva & sensitivity score & -7.17 & +1 & -4.48 \\
Wildenhain 2015 & YMR263W (SAP30) & hap & 3-octadecanamido-2 & z score & -81.84 & +1 & -10.87 \\
Hillenmeyer HOM & YBR058C (YBR058C) & dip & mitomycin c & z score & +56.29 & -1 & -15.95 \\
\bottomrule
\end{tabular}

\end{table}

Three things in that table are worth naming, because each is a place the datasets differ in a
way a reader would otherwise have to infer.

The Vanacloig strain perturbs four genes, not one. Three of them, PDR1, PDR3 and SNQ2, are
the drug-efflux regulators deleted to build the sensitized host, and they are the same three
in every Vanacloig record; the fourth is the gene being queried. That is what the four-gene
perturbation count in Table~\ref{tab:datasets} means, and it is the reason a Vanacloig
genotype contributes four indices to the perturbed set where the others contribute one.

One gene appears in four of the five rows on four different scales. UBP14 under crystal violet
reads $-3.40$ as a log ratio, the same gene under nocodazole reads $+1.92$ as a log ratio
pointing the other way, under an epothilone derivative $-7.17$ as a sensitivity score, and
under mitomycin C $+56.29$ as a z score. These are not phenotypes of different magnitude; they
are the same kind of measurement reported on scales that differ by a factor of twenty and, for
two of the five, with the opposite sign. The fifth row is a different gene because
Wildenhain's panel is 242 genes and does not contain this one, which is the narrowness
Table~\ref{tab:datasets} records.

After both steps all five land in the same range and the same direction. Every standardized
value in the table is negative, which under the shared convention means every one of these
records is a strain that grew badly, which is what picking the strongest response per dataset
was meant to produce.

\notesec{The two steps, and why in that order}
For a record from dataset $k$ with served response $y$,
\[
  \tilde y \;=\; o_k \cdot \frac{y - \mathrm{med}_k}{\mathrm{sd}_k},
\]
where $o_k \in \{-1, +1\}$ is the orientation factor and $\mathrm{med}_k$ and $\mathrm{sd}_k$
are the median and standard deviation of dataset $k$ on the TRAINING split only, so no
held-out value informs the scale. Orientation comes first because the median and the standard
deviation are computed on the oriented values; the standard deviation is unaffected by the
sign, but stating the order keeps the two from being applied inconsistently. Three of the five
sources already store a sick strain as negative and take $o_k = +1$, and only the two
Hillenmeyer arms flip.

Section~\ref{sec:joinability} is where $o_k$ comes from, and it is not a preference. The five
sources declare opposite conventions in their own words, and ignoring that puts two records of
the same gene under the same compound on opposite sides of zero.
